\scanpage{256}
\mspage{130}
% archive sheet number: 255

% slanted
\margin{en fait, ces conditions sont $=$ en l'absence de l'hypothèse Loc $(\unsure{6|5})$, cf 9.\unsure{9|8} ci-dessous}

\uline{Commentaire} \unsure{9.8}. Si les axiomes Loc $(1,2,3,\deleted{\unsure{6}}5)$\unsure{$^{*}$} sont satisfaits, alors l'axiome \inserted{de connexité} Loc $(8)$ \uline{équivaut} à chacune des suivantes :

a) $\mathrm{Hot}_W$ n'est équivalente ni à $e$, ni à $\Delta^1$

b) $\mathrm{Hot}_W$ \deleted{n'est pas une cat. \unsure{formée}}\inserted{a une infinité de classes} d'isomorphie

c) \deleted{$\unsure{e\to}$} $f_0:X_0\to Y_0$ (flèche donnée d'avance \unsure{dans|\deleted{dans}} $\mathrm{Cat}$, telle que $\pi_0(f_0)$ n'est pas iso, \inserted{et $X_0\ne\varnothing$}) n'est \uline{pas} un iso. de $\mathrm{Hot}_W$.

d) $X_0$ et $Y_0$ ne sont pas isomorphes dans $\mathrm{Hot}_W$ (où \unsure{$X.$|$X_0$}, \unsure{$Y.$|$Y_0$} sont deux objets donnés d'avance de $\mathrm{Cat}$, avec $X_0,Y_0$ non vides et $\mathrm{card}\,\pi_0(X)\ne\mathrm{card}\,\pi_0(Y)$)

% double vertical rule beside the opening four lines
\margin{\rule[-3\baselineskip]{0.4pt}{4\baselineskip}\hspace{1.4pt}\rule[-3.15\baselineskip]{0.4pt}{4.15\baselineskip}}

On peut donc dire que l'axiome de connexité joue le rôle d'un \uline{axiome de non trivialité} pour la cat. $\mathrm{Hot}_W$ localisée de $\mathrm{Cat}$. Il sera satisfait automatiquement pour tous les localiseurs \inserted{fond.} $W$ qu'on sera amené à construire ou à introduire. \deleted{Donc}\inserted{Mais}, on n'aura pratiquement jamais à l'invoquer comme une hypothèse, en vue de certaines conclusions.

$\backslash$

Je m'aperçois finalement que dans la prop. 9.7 l'hypothèse Loc $(6)$ est inutile :

% slanted
\margin{\unsure{Cf mieux une variante|Cf une version} 9.12 (p.~13\unsure{\deleted{3}\inserted{4}|3} bis)}

% slanted
\margin{Loc $(6\unsure{\deleted{a}|a})$}

\uline{Prop. 9.\unsure{9|8}}. Soit $W$ un localiseur fondamental de $\mathrm{Cat}$ (i.e. on a Loc $(1,2,3)$ sans plus), \inserted{\uline{plus} la condition que la \uline{somme} de $f_1,f_2\in W$ est \deleted{dans} $\in W$, $^{*}$} où l'axiome de connexité n'est pas satisfait, alors $\mathrm{Hot}_W$ est équivalent à la catégorie ponctuelle (si $\varnothing\to e$ est un iso dans $\mathrm{Hot}_W$)$^{**}$ ou à $\Delta^1$ (dans le cas contraire), et $\mathrm{Hot}^{\bullet}_W$ ($\overset{\mathrm{df}}{=}\mathrm{Hot}_W\smallsetminus\{\varnothing\}$) est équivalente à la cat. ponctuelle.)

% slanted
\margin{* Cette condition est vérifiée notamment si $W$ satisfait de plus l'un des trois axiomes Loc $(4)$ (cf 9.\unsure{1|2}) ; Loc $(5c')$ (cf 9.3) ou Loc $(6\deleted{\unsure{a|bis}})$.}

$(**)$ Cela signifie aussi qu'il existe $Z\ne\varnothing$ dans $\mathrm{Cat}$, tel que $\varnothing\to Z$ soit iso dans $\mathrm{Hot}_W$.
